\begin{lemma}[Four common missing colors for a PE hierarchy base]
Let $G$ be a finite loopless multigraph with distinguished edge $e_0$,
and let $\mathcal T$ be an extended Tashkinov tree whose generating
coloring is $\varphi$.  Put
\[
                    T_n=\mathcal T.\mathrm{tree}(n)
\]
at its terminal rung, and let $(T_{n,0}^{*},R_n,v_n)$ be a
\noderef{HierarchyBase} under $\varphi$.  Assume that the terminal
extension kind is PE, that $\varphi$ is proper and leaves exactly $e_0$
uncolored, that
\[
                         \Delta(G)+1\leq k,
\]
and that $V(T_{n,0}^{*})$ is elementary under $\varphi$.  Then
\[
             \bigl|\overline\varphi(T_n)
                 \cap\overline\varphi(R_n)\bigr|\geq4.
\]
\end{lemma}
\begin{proof}
Let $P=T_n(v_n)$ be the exact support prefix recorded by the hierarchy
base.  The PE clause of \noderef{HierarchyBase} says that the support
position is positive, that $R_n$ is a $\varphi$-Tashkinov closure of
$P$, and that $T_{n,0}^{*}$ is an
\noderef{OrderedTreeUnion} of $T_n$ and $R_n$.  In particular,
\[
                    V(T_{n,0}^{*})=V(T_n)\cup V(R_n).            \tag{1}
\]

We verify the endpoint hypotheses needed for the generic counting
lemma, rather than assuming them.  Let $j=0$ be the first edge index of
$T_n$.  The first-edge field of \noderef{TreeSequence} gives
$T_n.e_j=e_0$.  Its attachment field gives an earlier vertex index
$i\leq j$ such that the ends of this exact edge identity are the new
vertex $T_n.y_{j+1}$ and $T_n.y_i$.  Since $j=0$, necessarily $i=0$.
Thus, for
\[
                   x=T_n.y_1,\qquad y=T_n.y_0,
\]
one has
\[
                   \partial e_0=\{x,y\},qquad x,y\in V(T_n).   \tag{2}
\]

The support position is at least one.  Therefore both vertex indices
$0$ and $1$ occur in the exact prefix $P$, as defined by
\noderef{TreeSequencePrefix}; hence $x,y\in V(P)$.  The closure in the
base contains a
$\varphi$-\noderef{TashkinovAugmentationSequence} from $P$ to $R_n$, by
the definition of \noderef{TashkinovClosure}.
By \noderef{TashkinovAugmentationSequenceInitialSegment}, $P$ is an
exact ordered initial segment of $R_n$, so its vertex entries are retained.
Consequently
\[
                              x,y\in V(R_n).                     \tag{3}
\]

By (1), the assumed elementarity of $V(T_{n,0}^{*})$ is exactly the
elementarity of $V(T_n)\cup V(R_n)$.  Equations (2) and (3) show that
the two ends of the uniquely uncolored exact identity $e_0$ belong to
both constituent vertex sets.  All hypotheses of
\noderef{CommonMissingColorsAtUncoloredEndpoints} now hold: properness,
exactly-one-uncoloredness, and the degree bound are assumptions here,
while endpoint containment and union elementarity were just derived.
That lemma yields
\[
             4\leq
             \bigl|\overline\varphi(T_n)
                 \cap\overline\varphi(R_n)\bigr|,
\]
which is the desired assertion.
\end{proof}
